\section*{Chapitre \ref{ch:g10:freshwater-resources} --- Les ressources en eau douce}

\begin{solution}{exo:g10:freshwater-resources:1}
Un \omterm{def:g10:freshwater-resources:stock}{stock} est la quantité d'eau contenue dans un réservoir à un instant
donné, en \unit{km^3} (ou en kilogrammes). Un \omterm{def:g10:freshwater-resources:stock}{flux} est la quantité qui passe
d'un réservoir à un autre par unité de temps, en \unit{km^3} par an. Le
\omterm{def:g10:freshwater-resources:residence-time}{temps de résidence} est la durée moyenne que l'eau passe dans un réservoir,
en jours ou en années~; en régime permanent, c'est le \omterm{def:g10:freshwater-resources:stock}{stock} divisé par le
\omterm{def:g10:freshwater-resources:stock}{flux}.
\end{solution}

\begin{solution}{exo:g10:freshwater-resources:2}
Environ \qty{2.5}{\percent} de toute l'eau est douce. De l'eau douce,
\qty{68.7}{\percent} sont gelés dans les calottes, les glaciers et les
neiges permanentes, et \qty{30.1}{\percent} sont souterrains. Ce sont les
rivières qui en contiennent le moins, \num{2120}~\unit{km^3}.
\end{solution}

\begin{solution}{exo:g10:freshwater-resources:3}
L'\omterm{def:g10:freshwater-resources:renewable}{eau douce renouvelable} est l'eau douce renouvelée chaque année par les
précipitations, qui traverse une région en ruisselant et sous forme d'eaux
souterraines~: un volume par an, un \omterm{def:g10:freshwater-resources:stock}{flux}. L'\omterm{def:g10:freshwater-resources:fossil-groundwater}{eau souterraine fossile} est une
eau entrée dans un aquifère il y a longtemps et que la recharge actuelle ne
renouvelle pas~: c'est une réserve que l'on ne peut qu'entamer, un \omterm{def:g10:freshwater-resources:stock}{stock}.
\end{solution}

\begin{solution}{exo:g10:freshwater-resources:4}
$91\,000/2120 \approx 43$ fois plus.
\end{solution}

\begin{solution}{exo:g10:freshwater-resources:5}
L'atmosphère renouvelle son eau le plus vite (environ \num{9.5}~jours),
l'océan le plus lentement parmi les réservoirs calculés (environ
\num{3200}~ans)~; les eaux souterraines profondes sont plus anciennes
encore. $3200 \times 365.25/9.5 \approx \num{120000}$ fois plus longtemps.
\end{solution}

\begin{solution}{exo:g10:freshwater-resources:6}
Entrées~: $413\,000 + 73\,000 = \num{486000}$~\unit{km^3} par an, donc
$T = 12\,700/486\,000 \approx \num{0.0261}$ an $\approx \num{9.5}$ jours.
Sorties~: $373\,000 + 113\,000 = \num{486000}$~\unit{km^3}, autant~: l'air
est en régime permanent, et le \omterm{def:g10:freshwater-resources:residence-time}{temps de résidence} est le même.
\end{solution}

\begin{solution}{exo:g10:freshwater-resources:7}
$1\,335\,040/413 \approx \num{3230}$~ans. En régime permanent, l'océan
reçoit autant qu'il perd~: $373 + 40 = 413$ milliers de~\unit{km^3} par an,
si bien que le \omterm{def:g10:freshwater-resources:residence-time}{temps de résidence} calculé avec les entrées est le même.
\end{solution}

\begin{solution}{exo:g10:freshwater-resources:8}
Les terres reçoivent \num{113000}~\unit{km^3} par an~; \num{73000}
retournent à l'air par évaporation et transpiration et \num{40000}
ruissellent jusqu'à la mer~: $73\,000 + 40\,000 = \num{113000}$. La part
qui retourne à l'air vaut
$73/113 \approx \qty{65}{\percent}$.
\end{solution}

\begin{solution}{exo:g10:freshwater-resources:9}
$273.2 \times 10^6 \times 1233.5 \approx \qty{3.37e11}{m^3}$, soit
\qty{337}{km^3}~: presque autant que toute l'eau que contenait encore la mer
d'Aral en 1987.
\end{solution}

\begin{solution}{exo:g10:freshwater-resources:10}
$15.8 \times 0.3048 \approx \qty{4.82}{m}$. Sur les 65~ans qui vont de 1950
à 2015, $482/65 \approx \qty{7.4}{cm}$ par an.
\end{solution}

\begin{solution}{exo:g10:freshwater-resources:11}
$2/0.01 = 200$~ans. Le \omterm{def:g10:freshwater-resources:stock}{stock} baisse désormais de $0.05 - 0.01 =
\qty{0.04}{km^3}$ par an, si bien qu'il dure $2/0.04 = 50$~ans.
\end{solution}

\begin{solution}{exo:g10:freshwater-resources:12}
1960~: $1093/71.1 \approx 15$~ans~; 2010~: $98.1/11.4 \approx 8.6$~ans.
L'évaporation retire chaque année à peu près la même épaisseur d'eau de
chaque kilomètre carré du lac~; à mesure que le lac rétrécissait, sa
surface diminuait, et avec elle le volume total évaporé.
\end{solution}

\begin{solution}{exo:g10:freshwater-resources:13}
(1) La perte nette vaut $1.2 - 0.2 = \qty{1.0}{m}$ par an sur
$100 \times 10^6$~\unit{m^2}~: $10^8$~\unit{m^3}, soit \qty{0.1}{km^3} par an
doivent entrer. (2) L'apport tombe à \qty{0.05}{km^3}~: le lac perd
$\qty{0.05}{km^3} = 5 \times 10^7$~\unit{m^3}, une couche de
$5 \times 10^7/10^8 = \qty{0.5}{m}$. (3) À mesure que le niveau baisse, les
rives émergent et la surface du lac diminue~: moins de surface, moins
d'évaporation, jusqu'à ce que l'évaporation et l'apport s'équilibrent de
nouveau.
\end{solution}

\begin{solution}{exo:g10:freshwater-resources:14}
L'aquifère contient $100 \times 0.20 = \qty{20}{m}$ d'eau, soit
\qty{20000}{mm}~: à \qty{1}{mm} par an, \num{20000}~ans. L'eau profonde de
l'aquifère du Sahara est vieille de \num{200000} à un million d'années, dix
à cinquante fois plus~: elle a été rechargée sous des climats passés plus
humides.
\end{solution}

\begin{solution}{exo:g10:freshwater-resources:15}
$337/453\,000 \approx \qty{0.00074}{km}$, environ \qty{0.74}{m} d'eau. Le
niveau d'eau a baissé de \qty{4.8}{m}, environ six fois et demie plus.
L'eau n'occupe que les pores de la roche (et elle ne s'en écoule pas
entièrement)~: quand la surface de la nappe baisse d'un mètre, l'eau
libérée forme une couche bien plus mince qu'un mètre, ici environ
$0.74/4.8 \approx \num{0.15}$~m.
\end{solution}

\begin{solution}{pb:g10:freshwater-resources:1}
\textbf{1.} $1093/68\,000 \approx \qty{0.016}{km}$, environ \qty{16}{m}.

\textbf{2.} $1093 \times 10^9 \approx \qty{1.09e12}{m^3}$.

\textbf{3.} \qty{10}{g/L} valent \qty{10}{kg/m^3}~: $1.093 \times 10^{12} \times 10
\approx \qty{1.1e13}{kg}$, environ \num{11}~milliards de tonnes.

\textbf{4.} $1093/71.1 \approx 15$~ans.

\textbf{5.} Aucun fleuve n'en sort~: son eau ne peut partir que sous forme
de vapeur (ou en s'infiltrant dans le sol), tandis que le sel reste.

\textbf{6.} $71.1/68\,000 \approx \qty{0.00105}{km}$, environ \qty{1.05}{m}.

\textbf{7.} $9.4/68\,000 \approx \qty{0.00014}{km}$, environ \qty{0.14}{m}.

\textbf{8.} $71.1 - 9.4 = \qty{61.7}{km^3}$ par an.

\textbf{9.} $79 + 37 = \qty{116}{km^3}$~: $61.7/116 \approx \qty{53}{\percent}$,
environ la moitié.

\textbf{10.} Entrées $1.0 + 6.2 = \qty{7.2}{km^3}$~; sortie \qty{36.8}{km^3}~;
variation nette $7.2 - 36.8 = \qty{-29.6}{km^3}$.

\textbf{11.} $(502.7 - 327.2)/5 = \qty{35.1}{km^3}$ par an~: le même ordre
de grandeur que les \qty{29.6}{km^3} du bilan de 1987 (les années
diffèrent~: l'évaporation était de \qty{47.9}{km^3} en 1984).

\textbf{12.} $36.8/0.00105 \approx \qty{35000}{km^2}$.

\textbf{13.} $35\,000/68\,000 \approx \qty{52}{\percent}$, environ la moitié.

\textbf{14.} Le lac avait rétréci~: une surface moitié moindre évapore à peu
près moitié moins d'eau.

\textbf{15.} $345.6/1093 \approx \qty{32}{\percent}$, environ un tiers.

\textbf{16.} $10 \times 1093/345.6 \approx \qty{31.6}{g/L}$.

\textbf{17.} La salinité mesurée est plus faible~: une partie du sel avait
quitté l'eau, cristallisé sur le fond exondé et emporté par le vent sous
forme de poussières salées.

\textbf{18.} Environ \qty{1.05}{m} par an, l'épaisseur évaporée.

\textbf{19.} $7.2/35\,000 \approx \qty{0.0002}{km}$, environ \qty{0.20}{m}.

\textbf{20.} $1.05 - 0.20 \approx 0.85$ (\num{0.84} avec les valeurs non
arrondies)~: \textbf{en 1987, le niveau de la mer d'Aral a baissé
d'environ \qty{0.8}{m}, près d'un mètre en une seule année}.
\end{solution}
